\section{The vector correction on a regular ray}\label{app:coefficient}

The estimate of Proposition~\ref{prop:coupled} admits an explicit
leading coefficient when $a=\varepsilon$, $b=\kappa\varepsilon$
with fixed $0<\kappa<1$. Use the fixed Coulomb presentation of
\eqref{eq:fixedcoulomb} and write
\[
 \alpha_\varepsilon=\varepsilon u+O(\varepsilon^2),\qquad
 C=[u,\cdot],\quad
 G=(Q_0D_0^*D_0Q_0)^{-1},\quad
 \Lambda=\Pi_0C^*C\Pi_0.
\]
Here $u$ is the first real harmonic deformation; all remainders
hold in each fixed spatial regularity. The fixed inverse $G$
maps $Q_0L^2$ to $Q_0W^{2,2}$, with its higher elliptic versions.
No second derivative of the finite balancing correction is used.

The low spectral space is a graph over $K_0=\ker D_0$.
Let $X_\varepsilon:K_0\to Q_0L^2$ be its graph map and set
\begin{equation}\label{eq:graphcoefficient}
 R=-GQ_0C^*C\Pi_0,\qquad
 U_\varepsilon=(I+X_\varepsilon)
              (I+X_\varepsilon^*X_\varepsilon)^{-1/2}.
\end{equation}
Then
\begin{equation}\label{eq:graphexpansion}
 X_\varepsilon=\varepsilon^2R+O(\varepsilon^3),\qquad
 U_\varepsilon=I+\varepsilon^2R+O(\varepsilon^3).
\end{equation}
Indeed \eqref{eq:lowcomplement} gives $X_\varepsilon=O(\varepsilon^2)$.
Writing the blocks of $D_\varepsilon^*D_\varepsilon$ as
$A_{00},B,L_Q$, Coulomb orthogonality gives
$B=\varepsilon^2Q_0C^*C\Pi_0+O(\varepsilon^3)$ and
$A_{00}=\varepsilon^2\Lambda+O(\varepsilon^3)$.
The graph equation
$L_QX+B=X(A_{00}+B^*X)$, with
$L_Q=L_{Q,0}+O(\varepsilon)$ between fixed elliptic spaces,
proves \eqref{eq:graphexpansion}. Elliptic regularity upgrades
the finite-dimensional graph estimate to every required spatial norm.
The isometric normalization changes the parallel component only
at order $\varepsilon^4$.

Let $f_0$ be the zero-mean solution of \eqref{eq:Poisson} at
the enhancement point, and let $h=\Delta_0^{-1}f_0$ have zero
mean, with $\Delta_0=-\nabla^2$. Define
\begin{align}
 Hk&=\HH_2\bigl(\BF(Rk,F_0)+\BF(k,u\wedge u)\bigr),
 \label{eq:Hcoefficient}\\
 \langle k,Wl\rangle
   &=\int_S h\,\BF([u,k],[u,l])\,dV_0 .
 \label{eq:Wcoefficient}
\end{align}
Internal one-form indices in the second equation are contracted
with $g_0$. The harmonic transgression \eqref{eq:transgression}
and \eqref{eq:graphexpansion} give
\begin{equation}\label{eq:masscoefficient}
 U_\varepsilon^*M_\varepsilon U_\varepsilon
      =\varepsilon^4H^*H+O(\varepsilon^5).
\end{equation}
The terms linear in a second connection coefficient are exact
under harmonic projection, which is why no such coefficient
appears in \eqref{eq:Hcoefficient}.

For the kinetic term, $\Pi_0f_\varepsilon\Pi_0=0$ exactly.
Since $f_\varepsilon=f_0+O(\varepsilon)$,
\[
 U_\varepsilon^*f_\varepsilon U_\varepsilon
 =\varepsilon^2(\Pi_0f_0R+R^*f_0\Pi_0)+O(\varepsilon^3)
 =-2\varepsilon^2W+O(\varepsilon^3).
\]
For the last equality use $G(f_0k)=hk$ for parallel $k$:
$D_0^*D_0(hk)=(\Delta_0h)k=f_0k$. Self-adjointness of $G$
then gives
$\langle k,f_0Rl\rangle=-\int h\,\BF([u,k],[u,l])$.

There is a further cancellation in the energy-dependent term.
Parallelness implies $[F_0,k]=0$, so trace invariance gives
$I_0^*[u,k]=0$ pointwise. Thus
\[
 J_\varepsilon U_\varepsilon
 =I_\varepsilon^*D_\varepsilon U_\varepsilon
 =O(\varepsilon^2).
\]
The resolvent term in $Z_\varepsilon$ is consequently
$O(\varepsilon^4)$ on the fixed interval. Combining these
identities yields
\begin{equation}\label{eq:coefficientlimit}
 \varepsilon^{-4}U_\varepsilon^*
 E_\varepsilon(\varepsilon^2z)U_\varepsilon
 \longrightarrow H^*H+2zW
\end{equation}
uniformly for $z$ in compact sets. This is an operator statement
and remains valid at a repeated leading eigenvalue.

For the three simple positive branches let $k_j$ be the unit
$\Lambda$ eigenvectors with
$\ell_j=8(\kappa^2,1,1+\kappa^2)_j/(3V_0)$, and define
\begin{equation}\label{eq:sigmacoefficient}
 \sigma_j=\|Hk_j\|^2+
       2\ell_j\int_S h\,|[u,k_j]|^2\,dV_0.
\end{equation}
If $\mu_{\varepsilon,j}$ denotes the corresponding actual
leading gauge eigenvalue, the roots of \eqref{eq:pencil} obey
\begin{equation}\label{eq:raycorrection}
 \lambda_{\varepsilon,j}^{[1]}(\eta)-\mu_{\varepsilon,j}
 =\eta\varepsilon^4\sigma_j
 +O_\kappa(\eta\varepsilon^5+\eta^2\varepsilon^6).
\end{equation}
To justify the remainder, pull the normalized leading eigenvector
back by $U_\varepsilon$. Its distance from $k_j$ is
$O_\kappa(\varepsilon)$ and its spectral separation is at least
$c_\kappa\varepsilon^2$. Differentiating the simple pencil root
at $\eta=0$ gives
$\lambda'_j(0)=\langle p_j,E_\varepsilon(\mu_j)p_j\rangle$.
Equation~\eqref{eq:coefficientlimit} gives
$\lambda'_j(0)=\varepsilon^4\sigma_j+O_\kappa(\varepsilon^5)$.
On the low-root interval,
$E=O(\varepsilon^4)$, $E_\lambda=O(\varepsilon^2)$ and
$E_{\lambda\lambda}=O(\varepsilon^4)$. The separated eigenvector
equation therefore gives $\partial_\eta p_j=O_\kappa(\varepsilon^2)$
and $\partial_\eta^2\lambda_j=O_\kappa(\varepsilon^6)$.
Taylor's theorem in $\eta$ proves \eqref{eq:raycorrection}.

The Pauli family has an additional selection rule. Write
$u=U\otimes\sigma_3+\kappa V\otimes\sigma_1$, with $U,V$
fixed unitary adjoint-valued one-forms on $A$. For the generators
$\widetilde k_i=I_A\otimes i\sigma_i/\sqrt{6V_0}$,
\begin{equation}\label{eq:selectionrule}
 H\widetilde k_1=H\widetilde k_3=0.
\end{equation}
Products in $u\wedge u$ have multiplicity factors in
$\operatorname{span}\{I,\sigma_2\}$, so their pairings with
$\widetilde k_1,\widetilde k_3$ vanish. For these two inputs,
the double commutator $C^*C\widetilde k_i$ has multiplicity
factors only in $\operatorname{span}\{\sigma_1,\sigma_3\}$.
The fixed complementary projection and Green operator preserve
that subspace; pairing with $F_A\otimes I$ is again zero.
This proves both terms of \eqref{eq:Hcoefficient} vanish.
The smallest and middle branches thus have only the warping
integral in \eqref{eq:sigmacoefficient}.

The formula determines a background-dependent coefficient, not
a universal nonzero number. Its value requires the reference ASD
fields and normalized harmonic representatives. Zero mean does
not give the warping integral a sign. At $\kappa=1$, individual
branch coefficients require the geometric splitting on the
repeated leading eigenspace as well as \eqref{eq:coefficientlimit}.
Finally \eqref{eq:raycorrection} concerns the defined pencil:
its powers of $\eta$ do not bound uncomputed string corrections.
